\documentclass[../../../include/open-logic-section]{subfiles}

\begin{document}

\olfileid{sth}{cardinals}{cardsasords}
\olsection{作为序数的基数}

实际上，我们的基数理论将（毫不客气地）直接利用我们的序数理论。
也就是说：我们将直接把基数定义为某些特定的序数。
具体而言，我们给出如下定义：

\begin{defn}\ollabel{defcardinalasordinal}
如果 $A$ 可以良序化，那么 $\card{A}$ 就是使得 $\cardeq{A}{\gamma}$ 成立的最小序数 $\gamma$。
对于任意序数 $\gamma$，我们称 $\gamma$ 为\emph{基数}，当且仅当 $\gamma = \card{\gamma}$。
\end{defn}

我们刚才用了“$A$ 可以良序化”这个短语。正如数学中几乎总是如此，这里的模态表述不过是对一个存在性主张的粗略说法：说“$A$ \emph{可以}良序化”，就只是说“存在一个关系，它能良序化 $A$”。

但 \olref{defcardinalasordinal} 有一个问题。我们希望\emph{每个}集合都有一个大小，也就是说，对每个 $A$，$\card{A}$ 都存在。但我们刚刚给出的定义是以一个条件句开头的：“\emph{如果} $A$ 可以良序化……”。如果存在某个不能良序化的集合 $A$，那么我们的定义将根本无法定义对象 $\card{A}$。

所以，要使用 \olref{defcardinalasordinal}，我们需要一个保证：每个集合都能被良序化。遗憾的是，这个保证在 $\ZF$ 中是无法得到的。因此，如果我们想使用 \olref{defcardinalasordinal}，除了增加一条新公理之外别无他法，例如：

\begin{axiom}[良序公理]
每个集合都可以被良序化。
\end{axiom}

我们将在 \olref[choice][]{chap} 中讨论良序公理是否可接受。不过，从现在开始，我们将径自采用它。并且，利用它，可以非常直接地证明（按 \olref{defcardinalasordinal} 定义的）基数是存在的并且行为良好：

\begin{lem}\ollabel{lem:CardinalsExist}
对于任意集合 $A$：
\begin{enumerate}
	\item\ollabel{cardaexists} $\card{A}$ 存在且唯一；
	\item\ollabel{cardaapprox}  $\cardeq{\card{A}}{A}$；
	\item\ollabel{cardaidem}  $\card{A}$ 是一个基数，即
	$\card{A} = \card{\card{A}}$；
\end{enumerate}
\end{lem}

\begin{proof}
固定 $A$。由良序公理，存在一个良序结构 $\tuple{A, R}$。根据
\olref[ordinals][ordtype]{thmOrdinalRepresentation}，$\tuple{A,
 R}$ 同构于唯一的一个序数 $\beta$。所以
$\cardeq{A}{\beta}$。根据超限归纳法，存在唯一的最小序数 $\gamma$，使得 $\cardeq{A}{\gamma}$。所以 $\card{A}
= \gamma$，这就证明了 \olref{cardaexists} 和 \olref{cardaapprox}。
为证明 \olref{cardaidem}，由 $\gamma$ 的选取可知，若 $\delta \in \gamma$ 则 $\cardless{\delta}{A}$；又因等势是等价关系（\olref[sfr][siz][equ]{equinumerosityisequi}），故亦有 $\cardless{\delta}{\gamma}$。所以 $\gamma =
\card{\gamma}$。 
%So, for reductio, suppose that there is some ordinal $\delta \in \gamma$ such that $\cardeq{\gamma}{\delta}$. Then, $\cardeq{\cardeq{A}{\gamma}}{\delta}$ so that $\cardeq{A}{\delta}$ by \olref[sfr][set][equ]{equinumerosityisequi}, which contradicts the choice of $\gamma$ as the least ordinal such that $\cardeq{A}{\gamma}$. 
\end{proof}

下一个结果保证了 Cantor 原理乃至更多内容。
（注意基数继承了序数的序关系，即 $\cardfont{a} < \cardfont{b}$ 当且仅当 $\cardfont{a} \in \cardfont{b}$。
在以下表述中，我们将用哥特体字母表示已知为基数的对象。这相当标准。另一个常见的替代方案是用希腊字母——因为基数是序数——但选自字母表中部，例如：$\kappa, \lambda$。）：
\begin{lem}\ollabel{lem:CardinalsBehaveRight}
对于任意集合 $A$ 和 $B$：
\begin{align*}
	\cardeq{A}{B} &\text{ 当且仅当 } \card{A} = \card{B}\\
	\cardle{A}{B} &\text{ 当且仅当 } \card{A} \leq \card{B}\\
	\cardless{A}{B}&\text{ 当且仅当 } \card{A} < \card{B}
\end{align*}
\end{lem}

\begin{proof}
我们将证明第二个断言中从左到右的方向（其他情况类似，留作练习）。那么，考虑下图：
\begin{center}
	\begin{tikzpicture}
	\node (nodea) {$A$};
	\node[right = 6em of nodea] (nodeb) {$B$};
	\node[below = 2em of nodea] (nodecarda) {$\card{A}$};
	\node[below = 2em of nodeb] (nodecardb) {$\card{B}$};
	\draw[->] (nodea)--(nodeb);
	\draw[<->] (nodea)--(nodecarda);
	\draw[<->] (nodeb)--(nodecardb);
	\draw[->, dashed] (nodecarda)--(nodecardb);
\end{tikzpicture}
\end{center}
双箭头表示双射，其存在性由 \olref{lem:CardinalsExist} 保证。由 $\cardle{A}{B}$ 的假设，存在一个单射 $A\to B$。现在，从 $\card{A}$ 出发，经 $A$ 到 $B$ 再到 $\card{B}$ 沿箭头追踪，我们便得到了一个单射 $\card{A} \to \card{B}$（虚线箭头）。
\end{proof}\noindent 我们也可以利用 \olref{lem:CardinalsBehaveRight}
重新证明 Schr\"{o}der--Bernstein 定理。该定理断言：若
$\cardle{A}{B}$ 且 $\cardle{B}{A}$，则 $\cardeq{A}{B}$。我们曾将它陈述为
\olref[sfr][siz][sb]{thm:schroder-bernstein}，但最初是在
\olref[sfr][infinite][card-sb]{sec} 中——颇费周折地——证明的。
现在考虑：

\begin{proof}[施罗德-伯恩斯坦定理的重新证明]
若 $\cardle{A}{B}$ 且 $\cardle{B}{A}$，则由 \olref{lem:CardinalsBehaveRight} 得
$\card{A} \leq \card{B}$ 且 $\card{B} \leq \card{A}$。
于是由三分律及 \olref{lem:CardinalsBehaveRight} 得 $\card{A} = \card{B}$，
从而 $\cardeq{A}{B}$。
\end{proof}

\noindent
尽管这个证明非常简单，但它隐含地依赖于替换公理（用以确保
\olref[ordinals][ordtype]{thmOrdinalRepresentation}）和良序公理（用以保证
\olref{lem:CardinalsBehaveRight}）。相比之下，\olref[sfr][infinite][card-sb]{sec}
中的证明则自足得多（事实上，它可以在~$\Zminus$ 中完成）。

\end{document}